\problem{CM01}{Scheduling a Regular Bipartite Graph}{Hall's condition; matching decomposition}
There are $n$ workers and $n$ tasks. Each worker can perform exactly $d$ tasks, and each task can be performed by exactly $d$ workers, where $1\le d\le n$. Prove that $d$ rounds can be scheduled so that each worker performs a distinct task in every round and every allowed worker--task pair is used exactly once. Give the key proof of the matching theorem used.
\solution Draw the allowed pairs as a bipartite graph. For any worker set $S$, there are $d|S|$ incident edges, all ending in its neighbor set $N(S)$. Each task receives at most $d$ of these edges, so $d|S|\le d|N(S)|$, giving $|N(S)|\ge|S|$.
Hall's theorem states that this condition guarantees a matching covering all workers. Its proof is by induction on the left-side size. If a nonempty proper set $S$ has equality, match $S$ to $N(S)$ by induction. In the remaining graph, the original condition applied to $S\cup T$ gives $|N(T)\setminus N(S)|\ge|T|$ for every remaining worker set $T$, so induction also matches the complement.
If every nonempty proper set has strict inequality, choose an edge and remove its endpoints. Every remaining worker set loses at most one neighbor and still satisfies Hall's condition. Apply induction, then restore the chosen edge. A single vertex is the immediate base case.
The original graph thus has a perfect matching for the first round. Deleting its edges lowers every degree to $d-1$. Repeat on the remaining regular graph until degree zero, obtaining $d$ disjoint perfect matchings. Each round covers all workers and tasks, and all $nd$ allowed edges are used once.
\insight Degree counts turn a local regularity assumption into Hall's condition. Regularity survives removal of a matching, allowing a complete round-by-round decomposition.

\problem{CM02}{Congruences with Noncoprime Moduli}{Generalized CRT; extended Euclid}
For positive integers $m,n$ and integers $a,b$, determine when $x\equiv a\pmod m$ and $x\equiv b\pmod n$ are simultaneously solvable. Give an effective construction and the solution period. Then solve
\[
 x\equiv2\pmod6,\qquad x\equiv5\pmod9,
 \qquad x\equiv4\pmod5
\]
over the integers.
\solution Let $g=\gcd(m,n)$. Subtracting the congruences shows the necessary condition $g\mid b-a$. If it holds, write $x=a+mt$. The second congruence becomes
\[
 \frac m g t\equiv\frac{b-a}{g}\pmod{n/g}.
\]
The reduced coefficient and modulus are coprime. Extended Euclid computes the inverse of $m/g$ modulo $n/g$, and hence $t$. If $n/g=1$, there is no remaining restriction and one may take $t=0$. This proves sufficiency.
The difference of any two solutions is divisible by both $m$ and $n$, hence by $\operatorname{lcm}(m,n)=mn/g$. Conversely, adding any multiple of this number preserves both congruences. Thus the solutions form one residue class modulo $mn/g$.
In the example, put $x=2+6t$. Then $6t\equiv3\pmod9$, or $2t\equiv1\pmod3$, so $t\equiv2\pmod3$ and $x\equiv14\pmod{18}$. Write $x=14+18k$; the third congruence gives $3k\equiv0\pmod5$, hence $k\equiv0\pmod5$. All solutions are
\[
 \boxed{x=14+90j,\qquad j\in\mathbb Z}.
\]
\insight Noncoprime moduli require a compatibility check before inversion. Cancel the common divisor before seeking an inverse; do not divide by a noninvertible residue.

\problem{CM03}{Twelve Balls with Fixed Weighings}{Information bounds; ternary codes}
Exactly one of twelve identical-looking balls has abnormal mass, either heavier or lighter; all normal balls have equal mass. Using a balance without reference weights, with equal numbers of balls on the pans each time, design three weighings fixed in advance to identify the abnormal ball and whether it is heavy or light. Prove that two weighings are insufficient and that thirteen balls cannot be handled by this class of three predetermined weighings.
\solution Encode left-heavy, balance, and right-heavy as $1,0,-1$. There are $24$ ball-and-sign states, but only $3^2=9$ outcomes in two weighings, so at least three are required.
In the matrix below, columns represent balls and rows represent weighings: $1$ means left pan, $-1$ right pan, and $0$ omission.
\[
 \begin{array}{c|rrrrrrrrrrrr}
 &1&2&3&4&5&6&7&8&9&10&11&12\\\hline
 1&0&0&0&0&-1&-1&-1&1&1&1&-1&1\\
 2&0&-1&-1&-1&1&1&1&0&0&0&-1&1\\
 3&-1&1&0&-1&1&0&-1&-1&0&1&1&0
 \end{array}
\]
Each row has four positive and four negative entries, so all weighings are legal. Every column is nonzero, and no two are equal or opposite. Heavy ball $j$ produces column $j$ as its outcome vector; a light ball produces its negative. The $24$ possible outcomes are distinct, allowing unique decoding.
For thirteen balls, all $3^3-1=26$ nonzero outcomes would be needed. A column and its negative form a pair, so the thirteen columns would represent all such pairs. In any fixed coordinate, $2\cdot3^2=18$ of the $26$ nonzero vectors have a nonzero entry. Thus that row of the thirteen representative columns would contain exactly nine nonzero entries. Nine balls cannot be split equally between the pans, contradicting the balance constraint. Thirteen are impossible, and twelve achieves the maximum for this class of predetermined weighings.
\insight Counting outcomes gives only a necessary condition; a code must also satisfy physical balance constraints. Pairing a column with its negative handles the unknown heavy-or-light sign.
